\chapter{Perfect Secrecy and the One-Time Pad}

There is exactly one encryption scheme in this book that is provably, mathematically, unconditionally unbreakable. Not ``unbroken so far.'' Not ``would take a billion years.'' Unbreakable, full stop, against an adversary with infinite computing power and infinite time. It is the one-time pad, and Claude Shannon proved it secure in 1949.

That sounds like the end of cryptography. If we have a perfect cipher, why does the rest of this book exist? The answer is the most important lesson of this chapter, and it is a lesson about cost rather than about secrecy. The one-time pad's perfection comes with a price tag so steep that almost nobody can afford to pay it. Understanding precisely \emph{why} it is perfect, and precisely \emph{why} that perfection is unaffordable, is what pushes us off the cliff of perfect secrecy and into \emph{computational} security, where every other scheme in this book lives.

So this chapter does three things. It defines what ``perfectly secret'' even means, with enough care that we could hand the definition to a skeptic. It proves the one-time pad meets that definition. And it shows, with a theorem of Shannon's, that the pad's most annoying property -- a key as long as the message -- is not a flaw in the design but a law of nature for any perfectly secret scheme. After that, the two-time pad attack will show what happens the instant you try to cheat the law, and we will be ready to trade perfection for practicality.

\section{The setting: Alice, Bob, Eve, and a shared secret}
\label{ch09:setting}

Fix the cast. Alice wants to send a message $m$ to Bob over a channel that Eve can read in full. Alice and Bob share a secret key $k$ that Eve does not know. Alice computes a ciphertext $c = \Enc_k(m)$ and sends $c$; Bob computes $m = \Dec_k(c)$ and recovers the message. Eve sees $c$ and nothing else. This is \emph{symmetric} or \emph{private-key} encryption: one key, known to both legitimate parties, used to lock and to unlock.

\begin{definition}[Private-key encryption scheme]
\label{ch09:scheme-def}
A \emph{private-key encryption scheme} over message space $\mathcal{M}$, key space $\mathcal{K}$, and ciphertext space $\mathcal{C}$ is a triple of algorithms $(\Gen, \Enc, \Dec)$:
\begin{itemize}
  \item $\Gen$ outputs a key $k \in \mathcal{K}$, chosen according to some distribution (for us, almost always uniform).
  \item $\Enc$ takes a key $k$ and a message $m \in \mathcal{M}$ and outputs a ciphertext $c = \Enc_k(m) \in \mathcal{C}$. It may be randomized.
  \item $\Dec$ takes a key $k$ and a ciphertext $c$ and outputs a message $m = \Dec_k(c)$.
\end{itemize}
\emph{Correctness} requires that for every key $k$ that $\Gen$ can output and every message $m \in \mathcal{M}$,
\[
  \Dec_k(\Enc_k(m)) = m.
\]
\end{definition}

Correctness is the boring half: decryption undoes encryption for the legitimate user. Secrecy is the interesting half, and the rest of the chapter is about pinning down what it should mean.

Before we can talk about secrecy we have to be precise about what randomness is in play. There are two independent random things in this story, and conflating them is the source of most confusion.

The first is the key. We choose $k$ at random when we run $\Gen$. Write $K$ for the random variable that is the key.

The second is the message. From Eve's point of view, the message is also uncertain -- if she already knew $m$, there would be nothing to hide. We model her uncertainty with a random variable $M$ over $\mathcal{M}$, with some distribution $\Pr[M = m]$. In our model, this distribution is Eve's \emph{prior}: whatever she believed about the message before she saw any ciphertext. Maybe she knows Alice is voting in an election and $\mathcal{M} = \{\texttt{yes}, \texttt{no}\}$ with $\Pr[M=\texttt{yes}] = 0.7$. The scheme does not get to choose this distribution; the world does. A good definition of secrecy must work for \emph{every} possible prior.

\begin{focusbox}
Two random variables, kept rigorously separate:
\begin{itemize}
  \item $K$ -- the key, sampled by $\Gen$, unknown to Eve. The scheme controls its distribution.
  \item $M$ -- the message, modeling Eve's prior uncertainty. The world controls its distribution.
\end{itemize}
$K$ and $M$ are \textbf{independent}: Alice's choice of what to say has nothing to do with the random key her crypto library happened to generate. The ciphertext $C = \Enc_K(M)$ is then a random variable determined by both.
\end{focusbox}

\section{What should ``secret'' mean? Three tries}
\label{ch09:three-tries}

Let us reason our way to the definition rather than dropping it from the sky.

\textbf{Try 1: ``Eve cannot recover the key.''} Too weak, and aimed at the wrong target. We do not care about the key for its own sake; we care about the message. A scheme could leak the whole message while keeping the key safe -- imagine $\Enc_k(m) = m$, ignoring the key entirely. The key is perfectly hidden and the message is in plaintext. Useless.

\textbf{Try 2: ``Eve cannot recover the message.''} Better, but still too weak. Suppose Eve learns the first half of a message but not the second. She has not ``recovered the message,'' yet she has learned plenty. Or suppose the ciphertext reveals that $m$ is one of two possibilities and rules out a million others. Recovering nothing in full while learning almost everything is a real failure mode.

\textbf{Try 3: ``The ciphertext tells Eve nothing she did not already know.''} Now we are onto something. The right standard is that seeing $c$ should not change Eve's beliefs about the message at all. Whatever probability she assigned to each candidate message before, she assigns exactly the same probability after. The ciphertext is, informationally, a blank page.

That is Shannon's idea, and we can write it as a clean equation.

\begin{definition}[Perfect secrecy]
\label{ch09:perfect-def}
An encryption scheme $(\Gen,\Enc,\Dec)$ over message space $\mathcal{M}$ is \emph{perfectly secret} if for every distribution over $M$, every message $m \in \mathcal{M}$, and every ciphertext $c \in \mathcal{C}$ with $\Pr[C = c] > 0$,
\[
  \Pr[M = m \mid C = c] \;=\; \Pr[M = m].
\]
\end{definition}

Read it aloud: the probability the message was $m$, \emph{given that Eve saw ciphertext $c$}, equals the probability it was $m$ with no ciphertext at all. Conditioning on the ciphertext does nothing. In the language of the probability tools from earlier chapters, $M$ and $C$ are independent random variables.

\begin{intuitionbox}
Perfect secrecy is the statement ``$C$ and $M$ are independent.'' Eve can stare at the ciphertext forever, run it through any computation, cross-reference it with anything -- and her posterior belief about the message is identical to her prior. The ciphertext is correlated with the message through the key, but because the key is unknown and random, that correlation carries no usable information. The message and the ciphertext might as well live in separate universes.
\end{intuitionbox}

This is a strong demand. It must hold for \emph{every} prior distribution, so an adversary who happens to have excellent side information about the message is still no better off after seeing $c$. And nothing in the definition mentions computing power. Eve may be a deity with unbounded time. Perfect secrecy is an \emph{information-theoretic} guarantee: the information simply is not there to be extracted, regardless of how hard anyone computes.

\subsection{An equivalent and easier-to-check form}
\label{ch09:equiv}

Definition~\ref{ch09:perfect-def} quantifies over all priors, which is awkward to verify. Here is a reformulation that is logically equivalent and usually far easier to use.

\begin{lemma}[Ciphertext distribution is message-independent]
\label{ch09:equiv-lemma}
A scheme is perfectly secret if and only if for every pair of messages $m_0, m_1 \in \mathcal{M}$ and every ciphertext $c \in \mathcal{C}$,
\[
  \Pr[\Enc_K(m_0) = c] \;=\; \Pr[\Enc_K(m_1) = c],
\]
where the probability is over the choice of key $K \gets \Gen$ (and any randomness in $\Enc$).
\end{lemma}

In words: the distribution of the ciphertext is the same no matter which message you encrypt. Encrypting \texttt{ATTACK} and encrypting \texttt{RETREAT} produce ciphertexts drawn from \emph{identical} distributions. Since the ciphertext's distribution carries no fingerprint of the message, observing a particular $c$ cannot favor one message over another.

\begin{whybox}
\textbf{Why the two forms agree.} Start from the equivalent form: $\Pr[\Enc_K(m) = c]$ is some value $p(c)$ that does not depend on $m$. Apply Bayes' rule:
\[
  \Pr[M=m \mid C=c] = \frac{\Pr[C=c \mid M=m]\,\Pr[M=m]}{\Pr[C=c]}
  = \frac{p(c)\,\Pr[M=m]}{\Pr[C=c]}.
\]
Now $\Pr[C=c] = \sum_{m'} \Pr[C=c \mid M=m']\Pr[M=m'] = \sum_{m'} p(c)\Pr[M=m'] = p(c)$, because the probabilities of all messages sum to $1$. The $p(c)$ cancels top and bottom, leaving $\Pr[M=m\mid C=c] = \Pr[M=m]$. That is Definition~\ref{ch09:perfect-def}. The converse runs the same algebra backward. The whole proof is one application of Bayes plus the observation that a message-independent likelihood factors out of the normalizer.
\end{whybox}

We now have a target we can actually shoot at: build a scheme where every message produces the same spread of ciphertexts. The one-time pad does exactly this, and the mechanism is the humble XOR.

\section{XOR: the operation that makes pads work}
\label{ch09:xor}

The one-time pad runs on bitwise exclusive-or, written $\xor$. On single bits,
\[
  0 \xor 0 = 0,\quad 0 \xor 1 = 1,\quad 1 \xor 0 = 1,\quad 1 \xor 1 = 0.
\]
XOR is addition modulo $2$: it returns $1$ exactly when the inputs differ. On bit-strings of equal length we XOR position by position. Three properties make it the perfect locking mechanism.

\begin{itemize}
  \item \textbf{It is its own inverse.} $b \xor b = 0$ for any bit $b$, so $(x \xor k) \xor k = x$. XOR-ing with the same value twice returns you to the start. This single fact \emph{is} the decryption algorithm.
  \item \textbf{It is associative and commutative.} You can reorder and regroup XORs freely: $a \xor b \xor c$ is unambiguous.
  \item \textbf{XOR with a uniform random bit is a uniform random bit, and it hides the input.} This is the magic ingredient, so let us isolate it.
\end{itemize}

\begin{lemma}[The masking lemma]
\label{ch09:masking}
Let $x$ be any fixed bit (or bit-string) and let $K$ be a uniformly random bit (or bit-string of the same length), independent of $x$. Then $Y = x \xor K$ is uniformly distributed, and in particular its distribution does not depend on $x$.
\end{lemma}

\begin{proof}
Take the single-bit case; the multi-bit case is just this applied independently to each position. We want $\Pr[Y = 0]$ and $\Pr[Y = 1]$. If $x = 0$, then $Y = K$, so $\Pr[Y=0] = \Pr[K=0] = \tfrac12$ and $\Pr[Y=1] = \tfrac12$. If $x = 1$, then $Y = 1 \xor K$, which is $1$ when $K=0$ and $0$ when $K=1$; again each outcome has probability $\tfrac12$. Either way $Y$ is uniform. Since uniform is uniform regardless of $x$, the distribution of $Y$ does not depend on $x$.
\end{proof}

\begin{intuitionbox}
A fair coin XOR-ed onto your bit \emph{is} a fair coin. The mask $K$ is equally likely to leave your bit alone (when $K=0$) or to flip it (when $K=1$), and those two cases are equally likely, so the output is a coin flip no matter what you started with. Whatever you wanted to hide, the result looks like noise -- and looks like \emph{the same} noise whether your secret was a $0$ or a $1$. That last clause is Lemma~\ref{ch09:equiv-lemma} in miniature: the ciphertext's distribution is the same for every message. Stack $n$ independent fair coins side by side and you can hide an $n$-bit message under $n$ bits of perfect noise.
\end{intuitionbox}

This is the entire idea of the one-time pad. Everything below is bookkeeping on top of the masking lemma.

\section{The one-time pad}
\label{ch09:otp}

\begin{definition}[One-time pad]
\label{ch09:otp-def}
Fix a message length $n$. The message, key, and ciphertext spaces are all $\mathcal{M} = \mathcal{K} = \mathcal{C} = \{0,1\}^n$. The scheme is:
\begin{itemize}
  \item $\Gen$: output a key $k \getsr \{0,1\}^n$, i.e.\ $n$ uniformly random independent bits.
  \item $\Enc_k(m) = m \xor k$.
  \item $\Dec_k(c) = c \xor k$.
\end{itemize}
\end{definition}

Correctness is immediate from XOR being self-inverse:
\[
  \Dec_k(\Enc_k(m)) = (m \xor k) \xor k = m \xor (k \xor k) = m \xor 0^n = m.
\]

The data flow is one row of XOR gates with the key tapped in twice -- once to lock, once to unlock.

\begin{figure}[H]
\centering
\begin{tikzpicture}[
  font=\small,
  card/.style={draw=#1!70!black, fill=#1!10, rounded corners=3pt, line width=0.8pt, inner sep=4pt},
  gate/.style={draw=quotientorange!90!black, fill=quotientorange!18, circle, inner sep=1pt, minimum size=0.65cm, line width=0.9pt, font=\bfseries},
  arrowstyle/.style={-{Stealth[scale=1.0]}, line width=1.1pt}
]
  % Sender (Alice)
  \node[card=aliceblue, minimum width=2.4cm, minimum height=1.1cm, align=center] (m) at (0, 0)
    {\textbf{Alice (Sender)}\\[1pt]Plaintext $m \in \{0,1\}^n$};

  % Encryption XOR gate
  \node[gate] (x1) at (3.2, 0) {$\xor$};
  \node[card=structurepurple, minimum width=2.0cm, minimum height=0.85cm, align=center] (k1) at (3.2, -1.8)
    {Key $k \getsr \{0,1\}^n$\\[2pt]\scriptsize Shared secret pad};

  % Public Insecure Channel (Eve)
  \node[card=evered, minimum width=3.2cm, minimum height=1.1cm, align=center] (c) at (6.8, 0)
    {\textbf{Public Channel}\\[1pt]Ciphertext $c = m \xor k$};

  % Decryption XOR gate
  \node[gate] (x2) at (10.4, 0) {$\xor$};
  \node[card=structurepurple, minimum width=2.0cm, minimum height=0.85cm, align=center] (k2) at (10.4, -1.8)
    {Key $k \in \{0,1\}^n$\\[2pt]\scriptsize Identical pre-shared pad};

  % Receiver (Bob)
  \node[card=bobgreen, minimum width=2.4cm, minimum height=1.1cm, align=center] (m2) at (13.6, 0)
    {\textbf{Bob (Receiver)}\\[1pt]Recovered $m = c \xor k$};

  % Flows
  \draw[arrowstyle, aliceblue] (m.east) -- (x1.west);
  \draw[arrowstyle, structurepurple] (k1.north) -- (x1.south);
  \draw[arrowstyle, sourceblue] (x1.east) -- (c.west);

  \draw[arrowstyle, sourceblue] (c.east) -- (x2.west);
  \draw[arrowstyle, structurepurple] (k2.north) -- (x2.south);
  \draw[arrowstyle, bobgreen] (x2.east) -- (m2.west);

  % Eve observing
  \node[draw=evered!80, fill=evered!8, rounded corners=2pt, font=\scriptsize, align=center, inner sep=4pt] (eve) at (6.8, 1.5)
    {\textbf{Adversary Eve}\\[2pt]Sees $c$ only; posterior $\Pr[M=m \mid C=c] = \Pr[M=m]$};
  \draw[arrowstyle, evered, dashed] (c.north) -- (eve.south);

  % Secure transport line connecting keys
  \draw[dashed, structurepurple!80, line width=0.9pt] (k1.east) -- node[below=3pt, font=\scriptsize\bfseries, text=structurepurple, align=center] {Pre-shared via trusted courier\\[1pt]($n$ true random bits)} (k2.west);

\end{tikzpicture}
\caption{The one-time pad architecture. Alice and Bob pre-share an $n$-bit uniform random key $k$ through an out-of-band secure channel. Encryption bitwise XORs the plaintext with $k$, creating uniform noise on the wire. Bob recovers $m$ by XOR-ing $c$ with the same key, since $(m \xor k) \xor k = m \xor 0 = m$.}
\label{ch09:fig-otp-flow}
\end{figure}

\subsection{A tiny worked example}
\label{ch09:tiny}

Let us encrypt a 5-bit message by hand. Say Alice wants to send the message $m = \texttt{10110}$ and the key generator handed her $k = \texttt{01101}$. XOR them column by column -- output $1$ where the bits differ:
\[
\begin{array}{rccccc}
 m  = & 1 & 0 & 1 & 1 & 0 \\
 k  = & 0 & 1 & 1 & 0 & 1 \\ \hline
 c  = & 1 & 1 & 0 & 1 & 1
\end{array}
\]
So $c = \texttt{11011}$ goes on the wire. Bob, holding the same $k = \texttt{01101}$, XORs it back:
\[
\begin{array}{rccccc}
 c  = & 1 & 1 & 0 & 1 & 1 \\
 k  = & 0 & 1 & 1 & 0 & 1 \\ \hline
 m  = & 1 & 0 & 1 & 1 & 0
\end{array}
\]
and recovers \texttt{10110}. Notice that the third column shows $1 \xor 1 = 0$ in encryption and $0 \xor 1 = 1$ in decryption: the key bit toggles the message bit on the way out and toggles it back on the way in.

Here is the key property that makes this perfectly secret. Suppose Eve intercepts $c = \texttt{11011}$ and wonders whether the message was $\texttt{10110}$. The key that would produce $c$ from that message is $k = m \xor c = \texttt{10110} \xor \texttt{11011} = \texttt{01101}$ -- a perfectly possible key. But suppose she guesses the message was instead $\texttt{00000}$. The key explaining \emph{that} is $k' = \texttt{00000} \xor \texttt{11011} = \texttt{11011}$ -- also a perfectly possible key. In fact, \emph{for every one of the $2^5 = 32$ candidate messages, there is exactly one key that maps it to the observed $c$}, and all $32$ keys are equally likely. Eve cannot prefer any message over any other. The ciphertext has told her nothing.

\begin{walkbox}
You intercept the 3-bit ciphertext $c = \texttt{101}$. The message space is all of $\{0,1\}^3$. For the candidate message $m = \texttt{111}$, which key would Alice have needed? And for $m = \texttt{000}$? Convince yourself each of the $8$ candidate messages has its own unique explanatory key. \\[2pt]
\emph{Answer.} $k = m \xor c$. For $\texttt{111}$: $\texttt{111}\xor\texttt{101} = \texttt{010}$. For $\texttt{000}$: $\texttt{000}\xor\texttt{101}=\texttt{101}$. As $m$ ranges over all $8$ strings, $k = m \xor c$ also ranges over all $8$ strings, hitting each exactly once -- the map $m \mapsto m\xor c$ is a bijection. Eight messages, eight distinct equally likely keys, no winner.
\end{walkbox}

\subsection{The proof of perfect secrecy}
\label{ch09:otp-proof}

\begin{theorem}[Shannon, 1949]
\label{ch09:otp-secure}
The one-time pad is perfectly secret.
\end{theorem}

\begin{proof}
We verify the equivalent form, Lemma~\ref{ch09:equiv-lemma}: for every message $m$ and ciphertext $c$, the quantity $\Pr[\Enc_K(m) = c]$ is the same. Compute it. The encryption is $\Enc_K(m) = m \xor K$, so
\[
  \Pr[\Enc_K(m) = c] = \Pr[m \xor K = c] = \Pr[K = m \xor c].
\]
The last step XORs $m$ onto both sides of the event. Now $K$ is uniform on $\{0,1\}^n$, so every specific value -- including the value $m \xor c$ -- has probability exactly $2^{-n}$. Therefore
\[
  \Pr[\Enc_K(m) = c] = 2^{-n}
\]
for every $m$ and every $c$. The value $2^{-n}$ does not mention $m$ at all, so it is the same for all messages. By Lemma~\ref{ch09:equiv-lemma}, the scheme is perfectly secret.
\end{proof}

That is the whole proof. It is short because the masking lemma did the work: a uniform key turns the ciphertext into uniform noise whose distribution forgets the message. We can also see it through Definition~\ref{ch09:perfect-def} directly, which is worth doing once because it makes the ``Eve learns nothing'' story literal.

\begin{whybox}
\textbf{The Bayesian view.} Eve sees $c$ and updates her belief about $m$ using Bayes:
\[
  \Pr[M = m \mid C = c] = \frac{\Pr[C=c \mid M=m]\,\Pr[M=m]}{\Pr[C=c]}.
\]
The likelihood is $\Pr[C=c\mid M=m] = \Pr[K = m\xor c] = 2^{-n}$, a constant. The evidence term is $\Pr[C=c] = \sum_{m'} 2^{-n}\Pr[M=m'] = 2^{-n}$. The constants cancel:
\[
  \Pr[M=m\mid C=c] = \frac{2^{-n}\,\Pr[M=m]}{2^{-n}} = \Pr[M=m].
\]
Posterior equals prior. Whatever Eve believed before, she believes after. The ciphertext is, to her, a sealed box that updates nothing.
\end{whybox}

\begin{trapbox}
A frequent misreading: ``the one-time pad is secure because the key is secret.'' Secrecy of the key is necessary for every cipher, including breakable ones. What makes the pad \emph{perfect} is something sharper: the key is \emph{as long as the message} and \emph{uniformly random}, so the set of possible keys exactly covers the set of possible messages, one key per (message, ciphertext) pair. Shorten the key, reuse the key, or skew its distribution, and the bijection breaks -- some messages become more likely than others given $c$, and the proof collapses. The next sections show each of those three failures.
\end{trapbox}

\begin{historybox}
The construction is older than its proof. In 1882 the banker Frank Miller described a cipher using single-use random keys for telegraph traffic, but the idea was forgotten. In 1917 Gilbert Vernam, an engineer at AT\&T, patented an electromechanical cipher that XOR-ed teletype characters against a key tape; Joseph Mauborgne of the U.S. Army Signal Corps recognized that if the tape were truly random and never reused, the system was unbreakable. For decades that ``unbreakable'' was an engineering conviction, not a theorem.

It became a theorem in 1949 when Claude Shannon published \emph{Communication Theory of Secrecy Systems}, the paper that founded cryptography as a mathematical science. Shannon defined perfect secrecy, proved the pad achieves it, and -- the deeper result -- proved that \emph{any} perfectly secret scheme needs a key at least as long as the message (Section~\ref{ch09:key-length}). The same Shannon had, a year earlier, founded information theory itself. Cryptography and information theory are twins, born from one mind within twelve months.

The one-time pad saw heavy real use. Soviet intelligence used pads for its most sensitive cables; the Washington--Moscow ``hot line'' established after the Cuban Missile Crisis was protected by one-time tape. Where the key material can be distributed ahead of time by trusted courier and secrecy must be absolute, the pad still earns its keep.
\end{historybox}

\begin{implementbox}
A toy one-time pad in Python. It shows the mechanics -- and, by construction, every way to misuse it.
\begin{lstlisting}[language=Python]
import os

def otp_encrypt(message: bytes, key: bytes) -> bytes:
    if len(key) != len(message):
        raise ValueError("key must be exactly as long as the message")
    return bytes(m ^ k for m, k in zip(message, key))

# decryption is the SAME function: XOR is its own inverse
otp_decrypt = otp_encrypt

msg = b"ATTACK AT DAWN"
key = os.urandom(len(msg))          # fresh, uniform, message-length key
ct  = otp_encrypt(msg, key)
print(ct.hex())                      # looks like pure noise
assert otp_decrypt(ct, key) == msg   # round-trips perfectly

# The whole security argument lives in three rules this code does NOT enforce:
#   1. key from a TRUE random source (os.urandom is a CSPRNG, not perfect)
#   2. key length == message length (checked above)
#   3. NEVER reuse a key (nothing here stops you -- and that is the danger)
\end{lstlisting}
\textbf{Learning only.} This is a teaching toy. Real systems do not use one-time pads except in rare, specialized settings, precisely because the three rules above are nearly impossible to satisfy at scale. \texttt{os.urandom} is a cryptographic PRNG, not a source of perfect randomness, so even this code is only \emph{computationally} secure, not perfectly secret. Never roll your own cipher for production; use a vetted authenticated-encryption library.
\end{implementbox}

\section{Why the key must be long, random, and fresh}
\label{ch09:three-conditions}

The proof leaned on three properties of the key: it is uniformly random, it is as long as the message, and (implicitly, because we encrypt one message) it is used once. Drop any one and perfect secrecy dies. We take them in turn, with the most consequential -- reuse -- getting its own attack section.

\subsection{Random: a biased pad leaks}
\label{ch09:must-random}

The masking lemma needed $K$ \emph{uniform}. Suppose instead an implementation's key bits are biased, say each bit is $1$ with probability $0.9$. Then for a single-bit message, $\Pr[\Enc_K(0) = 1] = \Pr[K=1] = 0.9$ while $\Pr[\Enc_K(1) = 1] = \Pr[K=0] = 0.1$. The ciphertext bit's distribution now \emph{depends on the message}: a ciphertext $1$ screams ``the plaintext was probably $0$.'' Lemma~\ref{ch09:equiv-lemma} fails and so does secrecy. The keystream must be drawn from the uniform distribution, which in practice means a genuine entropy source, not a predictable counter or a low-entropy seed. The later chapter on randomness and CSPRNGs returns to how hard this is to get right.

\subsection{Long: a short key cannot cover the messages}
\label{ch09:must-long}

What if the key is shorter than the message -- say we have a $128$-bit key and want to send a $1$-megabyte file? The natural temptation is to stretch or repeat the key. Stretching it with an algorithm is exactly what a \emph{stream cipher} does (next chapter), and it gives up perfect secrecy for computational security. Repeating it is the classic Vigenere blunder and is catastrophically weak. But the deepest answer is Shannon's theorem: \emph{no} clever trick can make a short key perfectly secret over long messages. We prove that next.

\section{Shannon's theorem: perfect secrecy forces a huge key}
\label{ch09:key-length}

Here is the result that turns the one-time pad from a triumph into a cautionary tale. Its one impractical-looking property -- a key as long as everything you will ever send -- is not an artifact of XOR or of Vernam's tape. It is mandatory for \emph{any} perfectly secret scheme whatsoever.

\begin{theorem}[Shannon's key-length bound]
\label{ch09:shannon-bound}
If $(\Gen,\Enc,\Dec)$ is perfectly secret over message space $\mathcal{M}$ with key space $\mathcal{K}$, then $\abs{\mathcal{K}} \ge \abs{\mathcal{M}}$. In particular, the key must contain at least as many bits as the message.
\end{theorem}

\begin{proof}
The argument is a counting argument and is worth following slowly, because it explains the ``why'' viscerally. Assume for contradiction $\abs{\mathcal{K}} < \abs{\mathcal{M}}$.

Take any ciphertext $c$ that actually occurs, i.e.\ $\Pr[C = c] > 0$. Consider the set of messages you could possibly decrypt $c$ \emph{to}, ranging over all keys:
\[
  \mathcal{M}(c) = \set{\,\Dec_k(c) : k \in \mathcal{K}\,}.
\]
Each key $k$ yields one message $\Dec_k(c)$, so $\mathcal{M}(c)$ has at most $\abs{\mathcal{K}}$ elements -- you cannot reach more distinct messages than you have keys. Since we assumed $\abs{\mathcal{K}} < \abs{\mathcal{M}}$, the set $\mathcal{M}(c)$ misses at least one message. Call a missed message $m^{*}$: \emph{no key} decrypts $c$ to $m^{*}$.

But that means, if the true message had been $m^{*}$, the ciphertext $c$ could never have arisen: $\Pr[C = c \mid M = m^{*}] = 0$. Therefore, for an adversary who sees $c$,
\[
  \Pr[M = m^{*} \mid C = c] = 0.
\]
Now take a prior where $m^{*}$ has nonzero probability, say the uniform prior over $\mathcal{M}$, so $\Pr[M = m^{*}] > 0$. We have just shown $\Pr[M = m^{*} \mid C = c] = 0 \ne \Pr[M = m^{*}]$. The posterior differs from the prior, contradicting perfect secrecy (Definition~\ref{ch09:perfect-def}). Hence $\abs{\mathcal{K}} \ge \abs{\mathcal{M}}$.
\end{proof}

\begin{intuitionbox}
The picture behind the proof: perfect secrecy demands that \emph{every} message be a plausible explanation of \emph{every} ciphertext. Otherwise some ciphertext would rule a message out, and ruling out is information. To keep every message in play for a given $c$, you need at least one key per message -- a key that bends $c$ back to that message. With fewer keys than messages, some message has no key pointing to it from $c$, and the instant Eve sees $c$ she can cross that message off. Perfect secrecy cannot tolerate even a single crossed-off message. So: at least as many keys as messages.
\end{intuitionbox}

\begin{focusbox}
Read the consequence carefully. To send up to $2^n$ different messages with perfect secrecy you need at least $2^n$ keys, hence an $n$-bit key -- one fresh key bit for every message bit you will ever transmit. To protect a gigabyte you must have already shared a gigabyte of secret key, by some channel you trusted \emph{before} you needed the cipher. The cipher does not create secrecy; it \emph{transports} secrecy you already established. That is the whole catch, and it is unavoidable for any perfectly secret scheme.
\end{focusbox}

\begin{connectionbox}
This theorem is information theory wearing a cryptographic hat. In the language of entropy (companion probability book), perfect secrecy is equivalent to $H(M \mid C) = H(M)$ -- the ciphertext leaves the message's entropy untouched. Shannon's bound then follows from $H(K) \ge H(M)$: the key must carry at least as much uncertainty as the message it conceals. A short key has too little entropy to ``absorb'' a long message's uncertainty, and the leftover uncertainty has to show up somewhere -- namely, as information leaking into the ciphertext. The key-length bound and the entropy inequality are the same fact stated in two dialects.
\end{connectionbox}

\section{The two-time pad: reuse is fatal}
\label{ch09:two-time}

The word ``one'' in ``one-time pad'' is load-bearing. Reuse a single pad on two messages and the scheme does not degrade gracefully -- it shatters. This is the single most common way the pad is broken in the wild, and the attack is so clean it is worth working through bit by bit.

\begin{attackbox}
\textbf{Setup.} Alice encrypts two messages $m_1$ and $m_2$ with the \emph{same} key $k$ (the cardinal sin). Eve intercepts both ciphertexts:
\[
  c_1 = m_1 \xor k, \qquad c_2 = m_2 \xor k.
\]
\textbf{The break.} Eve XORs the two ciphertexts together. The shared key cancels:
\[
  c_1 \xor c_2 = (m_1 \xor k) \xor (m_2 \xor k) = m_1 \xor m_2 \xor (k \xor k) = m_1 \xor m_2.
\]
The key has vanished entirely (Figure~\ref{ch09:fig-twotime}). Eve now holds $m_1 \xor m_2$ -- the XOR of the two plaintexts, with no key involved. She has reduced an ``unbreakable'' cipher to the problem of separating two overlaid plaintexts, and that is a problem natural-language redundancy solves easily.

\textbf{Why $m_1 \xor m_2$ leaks the plaintexts.} English (and any structured data) is far from random. In ASCII, every letter has its top bits fixed; space is \texttt{0x20}; XOR-ing a space onto a letter flips the case bit, so $\text{letter} \xor \text{space}$ is just the same letter with flipped case. Eve does ``crib dragging'': she guesses a probable word like \texttt{"the~"} appears in $m_1$ at some position, XORs that guess into $c_1\xor c_2$ at that position, and reads off the corresponding slice of $m_2$. If the result is intelligible text, the guess was right and she has just recovered a chunk of \emph{both} messages at once. She slides the crib along, extends the readable region, and unzips the two messages letter by letter. With three or more reuses the job is trivial.
\end{attackbox}

\begin{example}[A reuse break in miniature]
\label{ch09:reuse-example}
Work with $6$-bit ``messages'' and watch only the XOR structure. Say
\[
  m_1 = \texttt{101100},\quad m_2 = \texttt{110010},\quad k = \texttt{011011}.
\]
Then
\[
  c_1 = m_1 \xor k = \texttt{110111},\qquad c_2 = m_2 \xor k = \texttt{101001}.
\]
Eve never learns $k$, but she computes
\[
  c_1 \xor c_2 = \texttt{110111} \xor \texttt{101001} = \texttt{011110}.
\]
Check against the plaintexts directly: $m_1 \xor m_2 = \texttt{101100} \xor \texttt{110010} = \texttt{011110}$. They match -- the key is gone. If Eve now learns or guesses even one of the messages, say she cribs $m_1 = \texttt{101100}$, she instantly gets the other: $m_2 = (c_1\xor c_2)\xor m_1 = \texttt{011110}\xor\texttt{101100} = \texttt{110010}$. One leaked message unlocks the partner for free, forever, with no key recovery needed.
\end{example}

\begin{figure}[H]
\centering
\begin{tikzpicture}[
  font=\small,
  card/.style={draw=#1!70!black, fill=#1!10, rounded corners=3pt, line width=0.8pt},
  gate/.style={draw=quotientorange!90!black, fill=quotientorange!18, circle, inner sep=1pt, minimum size=0.65cm, line width=0.9pt, font=\bfseries},
  arrowstyle/.style={-{Stealth[scale=1.0]}, line width=1.1pt}
]
  % Message 1 encryption
  \node[card=aliceblue, minimum width=2.2cm, minimum height=0.9cm] (m1) at (0, 1.2) {$m_1$};
  \node[gate] (x1) at (2.6, 1.2) {$\xor$};
  \node[card=structurepurple, minimum width=1.4cm, minimum height=0.7cm] (k) at (2.6, 0) {Key $k$};
  \node[card=evered, minimum width=2.4cm, minimum height=0.9cm] (c1) at (5.4, 1.2) {$c_1 = m_1 \xor k$};

  % Message 2 encryption
  \node[card=aliceblue, minimum width=2.2cm, minimum height=0.9cm] (m2) at (0, -1.2) {$m_2$};
  \node[gate] (x2) at (2.6, -1.2) {$\xor$};
  \node[card=evered, minimum width=2.4cm, minimum height=0.9cm] (c2) at (5.4, -1.2) {$c_2 = m_2 \xor k$};

  % Adversary XOR combination
  \node[gate] (x3) at (8.6, 0) {$\xor$};
  \node[card=targetgreen, minimum width=2.6cm, minimum height=1.0cm, align=center] (out) at (11.8, 0)
    {\textbf{Keyless XOR}\\[1pt]$m_1 \xor m_2$};

  % Flows
  \draw[arrowstyle, aliceblue] (m1.east) -- (x1.west);
  \draw[arrowstyle, structurepurple] (k.north) -- (x1.south);
  \draw[arrowstyle, evered] (x1.east) -- (c1.west);

  \draw[arrowstyle, aliceblue] (m2.east) -- (x2.west);
  \draw[arrowstyle, structurepurple] (k.south) -- (x2.north);
  \draw[arrowstyle, evered] (x2.east) -- (c2.west);

  \draw[arrowstyle, evered] (c1.east) -| (x3.north);
  \draw[arrowstyle, evered] (c2.east) -| (x3.south);
  \draw[arrowstyle, targetgreen] (x3.east) -- (out.west);

  % Annotations
  \node[font=\scriptsize\bfseries, text=trapred] at (2.6, 2.0) {Reuse of identical pad $k$!};
  \node[font=\scriptsize\bfseries, text=targetgreen!80!black, align=center] at (8.6, -1.6)
    {Key cancels completely:\\[1pt]$(m_1\xor k)\xor(m_2\xor k) = m_1\xor m_2$};

\end{tikzpicture}
\caption{The two-time pad disaster. When the same key $k$ is used to encrypt two distinct messages $m_1$ and $m_2$, Eve XORs the two ciphertexts together. Because $k \xor k = 0$, the secret key cancels entirely, leaving the raw bitwise sum $m_1 \xor m_2$ defenseless.}
\label{ch09:fig-twotime}
\end{figure}

\begin{figure}[H]
\centering
\begin{tikzpicture}[
  font=\small,
  cell/.style={draw=gray!50, fill=white, minimum width=0.68cm, minimum height=0.55cm, font=\ttfamily\footnotesize, inner sep=0pt, align=center},
  rescell/.style={draw=#1!70!black, fill=#1!15, minimum width=0.68cm, minimum height=0.55cm, font=\ttfamily\footnotesize\bfseries, inner sep=0pt, align=center},
  arrowstyle/.style={-{Stealth[scale=1.0]}, line width=1pt}
]
  % Header
  \node[font=\footnotesize\bfseries, text=structurepurple!90!black, anchor=west] at (0, 3.5)
    {Intercepted Ciphertext XOR Stream: $C_1 \xor C_2 = M_1 \xor M_2$ (15 bytes)};

  % Byte index ruler
  \foreach \i in {0,...,14} {
    \node[font=\tiny, text=gray!80] at (\i*0.78 + 0.39, 3.1) {\i};
  }

  % XOR stream byte cells (hex values)
  \def\xstream{{"1e","00","0b","10","00","19","0d","45","41","17","41","1d","07","41","40"}}
  \foreach \i in {0,...,14} {
    \pgfmathsetmacro{\val}{\xstream[\i]}
    \node[cell, fill=gray!8] (c\i) at (\i*0.78 + 0.39, 2.65) {\texttt{\val}};
  }

  % Case 1: Wrong Offset Trial (Offset 0)
  \begin{scope}[shift={(0, 0.75)}]
    \draw[draw=trapred!70, fill=trapred!3, rounded corners=3pt, dashed, line width=0.8pt] (-0.2, -0.6) rectangle (6.2, 1.25);
    \node[font=\scriptsize\bfseries, text=trapred!90!black, anchor=west] at (0, 1.0)
      {Trial A: Slide crib \texttt{"the "} at Offset 0};

    \node[font=\tiny\bfseries, text=gray!80, anchor=east] at (1.0, 0.5) {Crib $M_2$:};
    \foreach \i/\ch in {0/t, 1/h, 2/e, 3/\textvisiblespace} {
      \node[cell, fill=quotientorange!15] at (\i*0.78 + 1.45, 0.5) {\texttt{\ch}};
    }

    \node[font=\tiny\bfseries, text=gray!80, anchor=east] at (1.0, -0.15) {Result $M_1$:};
    \foreach \i/\ch in {0/j, 1/h, 2/n, 3/0} {
      \node[rescell=trapred] at (\i*0.78 + 1.45, -0.15) {\texttt{\ch}};
    }
    \node[font=\scriptsize\bfseries, text=trapred, anchor=west] at (4.7, 0.15)
      {\texttimes{} \shortstack[l]{\texttt{"jhn0"}\\Incoherent noise}};
  \end{scope}

  % Case 2: Correct Offset Trial (Offset 5)
  \begin{scope}[shift={(0, -1.45)}]
    \draw[draw=targetgreen!70, fill=targetgreen!3, rounded corners=3pt, line width=0.8pt] (3.8, -0.6) rectangle (11.8, 1.25);
    \node[font=\scriptsize\bfseries, text=targetgreen!90!black, anchor=west] at (4.0, 1.0)
      {Trial B: Slide crib \texttt{"the "} at Offset 5 (Match Found!)};

    \node[font=\tiny\bfseries, text=gray!80, anchor=east] at (5.0, 0.5) {Crib $M_2$:};
    \foreach \i/\ch in {0/t, 1/h, 2/e, 3/\textvisiblespace} {
      \node[cell, fill=quotientorange!15] at (\i*0.78 + 5.45, 0.5) {\texttt{\ch}};
    }

    \node[font=\tiny\bfseries, text=gray!80, anchor=east] at (5.0, -0.15) {Result $M_1$:};
    \foreach \i/\ch in {0/m, 1/e, 2/\textvisiblespace, 3/a} {
      \node[rescell=targetgreen] at (\i*0.78 + 5.45, -0.15) {\texttt{\ch}};
    }
    \node[font=\scriptsize\bfseries, text=targetgreen!90!black, anchor=west] at (8.7, 0.15)
      {\checkmark{} \shortstack[l]{\texttt{"me a"} \\ Valid English!}};
  \end{scope}

  % Downward cascade / unzipping explanation
  \node[draw=gray!30, fill=gray!2, rounded corners=3pt, font=\scriptsize, align=left, inner sep=5pt] at (5.9, -2.85)
    {\textbf{The Crib-Dragging Unzipping Cascade:}\\[2pt]
     1. Matching at offset 5 proves that $M_2[5..8]$ is \texttt{"the "} while $M_1[5..8]$ is \texttt{"me a"}.\\[1pt]
     2. English context readily suggests completing $M_1$ as \texttt{"meet me at noon"} and $M_2$ as \texttt{"send the cash.."}.\\[1pt]
     3. Eve plugs each extended word guess back into $C_1 \xor C_2$ to recover the other stream, unzipping both plaintexts completely.};

\end{tikzpicture}
\caption{Crib dragging against a two-time pad. Eve intercepts $C_1 \xor C_2 = M_1 \xor M_2$. She slides a common English crib (\texttt{"the "}) along the stream and XORs it with the bytes. At wrong offsets (Trial A), the result is gibberish like \texttt{"jhn0"}. At the correct offset 5 (Trial B), valid English plaintext (\texttt{"me a"}) emerges in $M_1$, confirming the guess and simultaneously unlocking both plaintexts.}
\label{ch09:fig-crib-dragging}
\end{figure}

\begin{historybox}
\textbf{VENONA.} During and after World War II, Soviet agents encrypted cables with one-time pads. Under wartime pressure, the pad-printing operation duplicated some pages of key material -- the same pad used on more than one message. U.S. and British cryptanalysts (the VENONA project, 1943--1980) detected the reuse and exploited exactly the $c_1 \xor c_2$ collapse above, painstakingly separating overlaid plaintexts. The decrypts exposed espionage networks, including the leak of atomic-bomb secrets. A theoretically perfect cipher was defeated not by breaking the math but by an operational shortcut that violated the ``one-time'' rule. The lesson is the recurring lesson of this whole book: schemes fail at the seams, where assumptions meet reality.
\end{historybox}

\begin{attackbox}
\textbf{Modern echo: WEP and stream-cipher reuse.} The same break sank Wi-Fi's first security protocol. WEP used the RC4 stream cipher -- effectively a pad generated from a key plus a $24$-bit ``initialization vector'' (IV). The IV is sent in the clear and is supposed to make each packet's pad unique. But $24$ bits is tiny: on a busy network the IV cycles in hours, and by the birthday bound (companion probability book) collisions appear much sooner. Two packets with the same IV reuse the same pad, and an eavesdropper recovers $p_1 \xor p_2$ exactly as above. Known plaintext structure in network packets (fixed headers, predictable handshake bytes) gives the cribs for free. WEP is the two-time pad attack, mass-produced -- and it is why WEP was abandoned. The repair, carried through the rest of this book, is the \emph{nonce}: a per-message value guaranteed never to repeat, so a key plus nonce never produces the same keystream twice.
\end{attackbox}

\begin{trapbox}
``I will just hash the key differently each time'' or ``I will append a counter to the message'' are not fixes -- they are new ciphers that need their own security analysis. With a genuine one-time pad there is no patch for reuse: the only correct behavior is to advance to fresh, never-before-used key material for every message, exactly $\abs{m}$ bits of it. The moment you find yourself wanting to reuse pad bytes ``just this once,'' you have outgrown the one-time pad and need a stream or block cipher with a nonce.
\end{trapbox}

\section{Stepping off the cliff: why we settle for computational security}
\label{ch09:computational}

Add up the bill. Shannon's bound (Theorem~\ref{ch09:shannon-bound}) says a perfectly secret scheme needs a key as long as the total volume of traffic it will ever protect. That key must be uniformly random, transported to both parties over an already-secure channel before any communication, kept perfectly secret, and \emph{never} reused under penalty of total collapse (Section~\ref{ch09:two-time}). For two people sharing a suitcase of one-time tape ahead of a mission, fine. For your browser opening a TLS connection to a server it has never met, having shared no prior secret, intending to exchange megabytes -- impossible. You cannot pre-share a gigabyte of perfect key with every website on Earth.

So we make a trade that defines all of modern cryptography. We give up \emph{perfect} secrecy and ask instead for \emph{computational} secrecy: the ciphertext may, in principle, contain information about the message, but extracting it requires solving a problem so hard that no realistic adversary -- bounded to, say, $2^{128}$ operations over the lifetime of the universe -- can do it.

\begin{focusbox}
The two pillars of the relaxation, in contrast to the perfect-secrecy bullets above:
\begin{itemize}
  \item \textbf{Bounded adversary.} We stop worrying about a computationally unlimited Eve and only defend against a probabilistic polynomial-time ($\PPT$) one. A deity could break our ciphers; no machine that will ever be built can.
  \item \textbf{Negligible advantage.} We allow the adversary a \emph{tiny} probability of success -- but tiny in the strong sense of \emph{negligible} ($\negl$): smaller than $1/p(n)$ for every polynomial $p$, as the security parameter $n$ grows. Doubling the key size makes the break time explode while the adversary's success chance vanishes.
\end{itemize}
In exchange for these two concessions we win everything the pad denied us: short keys (a $256$-bit key protecting terabytes), key reuse under proper nonce discipline, and -- once we add public-key methods -- secure communication with strangers we never pre-shared a secret with.
\end{focusbox}

\begin{whybox}
\textbf{Is the trade safe?} It is conditional. Perfect secrecy needs no assumptions -- it is a theorem. Computational secrecy rests on \emph{unproven} hardness assumptions: that factoring large integers is slow, that the AES function behaves like a random permutation, that certain discrete logarithms are intractable. If one of those assumptions falls, the schemes built on it fall too. We accept this because the alternative -- perfect secrecy -- is unusable at scale, and because the assumptions have survived decades of the world's best cryptanalysts trying to break them. The companion algorithms book quantifies exactly what ``$2^{128}$ operations'' buys you and why doubling a key length squares the attacker's work. The rest of \emph{this} book is, in a sense, the engineering of living comfortably with that trade.
\end{whybox}

\begin{applicationbox}
\textbf{Where the pad survives, and where computational security takes over.}
\begin{itemize}
  \item \textbf{One-time pad, still used:} diplomatic and intelligence links where couriers can pre-distribute key tape and the secrecy requirement is absolute and indefinite. Numbers stations -- shortwave broadcasts of spoken digit groups -- are widely understood to deliver one-time-pad key material to field agents. Quantum key distribution (QKD) is, in effect, a way to grow fresh shared random bits for use as a pad.
  \item \textbf{Computational security, everywhere else:} TLS (the \texttt{https} lock), the WPA2/WPA3 that replaced broken WEP, Signal's end-to-end messaging, SSH, disk encryption, your phone. None of these pre-share a message-length key; all of them rely on bounded adversaries and unproven hardness. They give up Shannon's perfection and in return scale to the entire Internet.
\end{itemize}
The mental model to carry forward: the one-time pad is the gold standard nobody can afford, and everything that follows is a carefully engineered counterfeit -- indistinguishable from gold to any buyer who cannot spend more than $2^{128}$ operations checking.
\end{applicationbox}

\section{What to remember}
\label{ch09:summary}

Perfect secrecy means the ciphertext is statistically independent of the message: posterior equals prior, for every prior, against an adversary of unlimited power (Definition~\ref{ch09:perfect-def}). The one-time pad achieves it because a uniform key masks any message into uniform noise whose distribution forgets which message it was (Theorem~\ref{ch09:otp-secure}), and the masking lemma (Lemma~\ref{ch09:masking}) is the one gear the whole machine turns on.

The price is set by Shannon's bound (Theorem~\ref{ch09:shannon-bound}): the key must be at least as long as the message, for \emph{any} perfectly secret scheme, no exceptions, no clever encodings. The key must also be uniformly random and used exactly once -- reuse collapses two ciphertexts to $m_1 \xor m_2$ and hands the plaintexts to anyone (Section~\ref{ch09:two-time}), as VENONA and WEP learned the hard way.

Because that price is unpayable for Internet-scale communication, we descend from perfect secrecy to computational secrecy: bounded adversaries, negligible advantage, short reusable keys, security resting on hardness assumptions instead of theorems. That descent is the doorway to every other chapter in this book.

\section*{Exercises}
\addcontentsline{toc}{section}{Exercises}

\begin{exercise}[XOR warm-up]
\label{ch09:ex-xor}
Let $m = \texttt{11010010}$ and $k = \texttt{01100101}$. Compute $c = m \xor k$. Then verify $c \xor k = m$. Finally, suppose Eve learns that the ciphertext is $c$ and the message is $m$; show she can compute the key, and explain in one sentence why this does not violate perfect secrecy.
\end{exercise}

\begin{exercise}[Every message is reachable]
\label{ch09:ex-reachable}
Fix a one-time-pad ciphertext $c \in \{0,1\}^n$. Prove that the map $m \mapsto m \xor c$ is a bijection on $\{0,1\}^n$, and explain how this bijection is exactly the statement that every candidate message has exactly one explanatory key. Relate this to the counting step in the proof of Shannon's bound (Theorem~\ref{ch09:shannon-bound}).
\end{exercise}

\begin{exercise}[Posterior equals prior, by hand]
\label{ch09:ex-bayes}
Let $\mathcal{M} = \{0,1\}^2$ and suppose Eve's prior is $\Pr[M=\texttt{00}]=\tfrac12$, $\Pr[M=\texttt{01}]=\tfrac14$, $\Pr[M=\texttt{10}]=\tfrac18$, $\Pr[M=\texttt{11}]=\tfrac18$. Using a one-time pad with a uniform $2$-bit key, compute $\Pr[M=\texttt{00}\mid C=\texttt{10}]$ from first principles with Bayes' rule, and confirm it equals the prior $\tfrac12$.
\end{exercise}

\begin{exercise}[Break a two-time pad]
\label{ch09:ex-twotime}
You intercept two ciphertexts encrypted under the \emph{same} 8-bit pad:
\[
  c_1 = \texttt{10011100}, \qquad c_2 = \texttt{11001010}.
\]
Compute $c_1 \xor c_2$. You are then told (a crib) that $m_1 = \texttt{01000001}$. Recover $m_2$ without ever learning the key, and state the general formula you used.
\end{exercise}

\begin{exercise}[The bias leak]
\label{ch09:ex-bias}
An implementation generates each key bit independently as $1$ with probability $p \ne \tfrac12$. For a one-bit message, compute $\Pr[M=0 \mid C=1]$ under a uniform prior $\Pr[M=0]=\Pr[M=1]=\tfrac12$, as a function of $p$. Show it equals $\tfrac12$ only when $p=\tfrac12$, and interpret what a deviation tells the adversary.
\end{exercise}

\begin{exercise}[A scheme that is not perfectly secret]
\label{ch09:ex-notperfect}
Define a ``shift cipher'' over $\mathcal{M}=\Zn{26}$: $\Gen$ picks $k \getsr \Zn{26}$, $\Enc_k(m) = (m+k) \bmod 26$. Prove that for a \emph{single} letter this scheme \emph{is} perfectly secret. Then argue that for two-letter messages drawn from English (so $M$ is non-uniform), encrypting both letters with the same $k$ is no longer perfectly secret, and exhibit a ciphertext-message pair whose posterior differs from its prior.
\end{exercise}

\begin{exercise}[Shannon's bound is tight]
\label{ch09:ex-tight}
The one-time pad meets Shannon's bound with equality: $\abs{\mathcal{K}} = \abs{\mathcal{M}} = 2^n$. Construct a \emph{different} perfectly secret scheme over $\mathcal{M} = \Zn{q}$ (any modulus $q$, not just a power of two) with $\abs{\mathcal{K}} = \abs{\mathcal{M}} = q$, prove it perfectly secret, and confirm it too meets the bound with equality. (Hint: replace XOR by addition modulo $q$.)
\end{exercise}

\begin{exercise}[Why pre-sharing is the catch]
\label{ch09:ex-presharing}
Estimate the amount of one-time-pad key material two people would need to pre-share to protect one year of typical messaging traffic (pick and state a reasonable bytes-per-day figure). Compare this to the $32$ bytes of a typical symmetric key in a computationally secure scheme. Use the contrast to explain, in a short paragraph, why TLS cannot use a one-time pad.
\end{exercise}

\begin{exercise}[Perfect secrecy via entropy]
\label{ch09:ex-entropy}
Using the entropy formulation from the companion probability book, perfect secrecy is equivalent to $H(M \mid C) = H(M)$. Assuming this, prove $H(K) \ge H(M)$, and explain how this recovers Shannon's key-length bound. (Hint: the message is a deterministic function of the key and ciphertext together, so $H(M \mid C) \le H(K \mid C) \le H(K)$.)
\end{exercise}

\begin{exercise}[Malleability of the pad]
\label{ch09:ex-malleable}
The one-time pad is perfectly \emph{secret} but offers no \emph{integrity}. Suppose Eve knows that $m$ is an account number and that the ciphertext is $c = m \xor k$. Show that by XOR-ing a chosen mask $\delta$ into $c$, Eve can flip specific bits of the decrypted message to a value of her choosing, \emph{without knowing $k$}. Give a concrete $4$-bit example, and explain why this means perfect secrecy does not imply security against active tampering -- motivating the message authentication of a later chapter.
\end{exercise}
